\subsection*{一步一步建立直观}

\paragraph{1. 把损失想成脚下的地形}

右图采用 $L(w)=w^2$，最低点在 $w=0$。导数告诉我们在当前位置向右走会使损失增加还是减少，所以更新时沿导数的反方向迈步。步长还乘以学习率；它既不是导数本身，也不是每次必须走固定距离。



\paragraph{2. 跨过最低点不一定表示失败}

把学习率提高后，轨迹可能左右交替。写出误差递推 $w_{k+1}=(1-2\eta)w_k$，就能看出只要这个倍率的绝对值小于一，交替仍会逐步缩小。图形里的来回摆动需要结合幅度判断，不能一见振荡就断言不收敛。



\paragraph{3. 过大的更新会放大误差}

再提高学习率，误差在两侧反复放大。这里目标函数光滑而且只有一个最低点，算法仍然会失败，说明步长本身是关键条件。高维神经网络的曲率随位置变化，不能把这个玩具函数的稳定区间直接照搬成通用学习率建议。



\subsection*{定义、命题与证明}

\paragraph{定义：梯度下降的一维更新}

对可微损失 $L$，给定学习率 $\eta$，按 $w_{k+1}=w_k-\eta L'(w_k)$ 更新参数。学习率控制使用局部斜率的步幅。

\paragraph{命题：二次损失的精确稳定区间}

对 $L(w)=w^2$，任意初值的迭代在 $0<\eta<1$ 时收敛到零；非零初值在 $\eta>1$ 时误差绝对值增长。

\begin{proof}

因为 $L'(w)=2w$，更新式成为 $w_{k+1}=(1-2\eta)w_k$，从而 $w_k=(1-2\eta)^kw_0$。

几何序列趋零当且仅当倍率的绝对值小于一（非零初值）。解 $|1-2\eta|<1$ 得 $0<\eta<1$。端点 $\eta=0$ 不更新，$\eta=1$ 对非零初值等幅往返，均不属于该收敛区间。

\end{proof}
